\chapter{Finite Fields and Arithmetic Frobenius}
\label{ch:finite-fields}

Let \(k\) be a finite field of cardinality \(q=p^f\), with \(p\)
prime, and fix an algebraic closure \(\Omega\) of \(k\).
The finite subfields used here are literal subfields of this one
closure. This choice makes all inclusion and restriction maps part
of the same system.

\section{Finite layers of the algebraic closure}

\begin{proposition}[Root fields in the chosen closure]\label{res:E01}
For each \(n\geq 1\), the roots of \(X^{q^n}-X\) in \(\Omega\)
form a subfield \(k_n\) with \(q^n\) elements. This field contains
\(k\), has \([k_n:k]=n\), and is a finite Galois extension of \(k\).
It is the unique subfield of \(\Omega\) with that cardinality.
\end{proposition}
\begin{proof}
The polynomial splits in \(\Omega\) and its derivative is \(-1\).
All its \(q^n\) roots are therefore distinct. The map
\(a\mapsto a^{q^n}\) is an iterate of the characteristic-\(p\)
Frobenius endomorphism. Its fixed elements contain \(0\) and \(1\)
and are closed under addition, subtraction, multiplication and
inversion of nonzero elements. They form a field. Every \(a\in k\)
satisfies \(a^q=a\), by the order of \(k^\times\) for \(a\ne0\),
so \(k\subseteq k_n\). The cardinality of a finite-dimensional
\(k\)-space of dimension \(d\) is \(q^d\), giving \([k_n:k]=n\).
Since \(k_n\) consists of all roots, it is the splitting field of
this separable polynomial. Finally, in a field of cardinality
\(q^n\) every element satisfies \(a^{q^n}=a\), again by the order of
the multiplicative group. Any such subfield of \(\Omega\) is therefore
contained in \(k_n\), and equality follows by counting.
Compare \cite[Propositions 4.20 and 4.23]{milne},
\cite[Theorem 2.5]{ln}, and \cite[Chapter 13]{d07}.
\end{proof}

\begin{proposition}[The divisor lattice]\label{res:E02}
For positive integers \(m,n\),
\[
\begin{aligned}
   k_m\subseteq k_n &\quad\Longleftrightarrow\quad m\mid n,\\
   k_m\cap k_n &= k_{\gcd(m,n)},&
   k_m k_n &= k_{\operatorname{lcm}(m,n)} .
\end{aligned}
\]
Here \(k_m k_n\) denotes the compositum in \(\Omega\).
The inclusions for \(m\mid n\) are \(k\)-embeddings and satisfy
the identity and composition laws.
\end{proposition}
\begin{proof}
If \(m\mid n\), an element fixed by the \(q^m\)-power map is fixed by
its \(n/m\)-th iterate. In the reverse direction, use the tower formula
for degrees and Proposition~\ref{res:E01}.
The degree of the intersection divides both \(m\) and \(n\). The
intersection contains \(k_{\gcd(m,n)}\), so its degree is their gcd.
The compositum lies in \(k_{\operatorname{lcm}(m,n)}\). Its degree
is a multiple of both \(m\) and \(n\), so equality holds. The embeddings
are inclusions of subfields of one field. Applying two successive
inclusions to an element leaves that element unchanged, which gives
the composition law. See \cite[Corollary 4.21 and Proposition 4.23]{milne}
and \cite[Theorem 2.6]{ln}.
\end{proof}

\begin{corollary}[Exhaustion by finite fields]\label{res:E03}
The \(k_n\) are all the finite intermediate fields of \(\Omega/k\),
and
\[
                \Omega=\bigcup_{n\geq1}k_n=\bigcup_{j\geq1} k_{j!}.
\]
The factorial stages are cofinal in the divisibility order.
\end{corollary}
\begin{proof}
A finite intermediate field of degree \(d\) has \(q^d\) elements, so
by Proposition~\ref{res:E01} it is \(k_d\). For \(a\in\Omega\) the
simple field \(k(a)\) is a finite intermediate field, since \(a\)
is algebraic over \(k\). This proves the first union.
Since \(n\mid n!\), Proposition~\ref{res:E02} proves the second union
and cofinality. The factorial stages are Galois by
Proposition~\ref{res:E01}. Their directed union is Galois by
Theorem~\ref{res:F04}, and its finite Galois subextensions are
exactly the fields in this divisibility system.
See \cite[Propositions 4.23--4.24]{milne}.
\end{proof}

\section{Frobenius as a compatible generator}

\begin{theorem}[Arithmetic Frobenius and its restrictions]\label{res:E04}
The map \(\Phi:\Omega\to\Omega\), \(x\mapsto x^q\), is a
\(k\)-automorphism of infinite order. Write
\(\Phi_n=\Phi|_{k_n}\). Then \(\Phi_n\) has order \(n\) and
generates \(\operatorname{Gal}(k_n/k)\). For \(m\mid n\) the
group \(\operatorname{Gal}(k_n/k_m)\) is generated by
\(\Phi_n^m\) and has order \(n/m\).
\end{theorem}
\begin{proof}
The characteristic-\(p\) Frobenius and hence its \(f\)-th iterate
are field endomorphisms. The latter is injective and fixes \(k\).
It preserves every root field \(k_n\), and on that finite set is
bijective. By Corollary~\ref{res:E03} it is surjective on \(\Omega\).

By construction \(\Phi_n^n=1\). If \(0<d<n\) and
\(\Phi_n^d=1\), every element of \(k_n\) would be a root of
\(X^{q^d}-X\). Its degree is smaller than \(|k_n|\), which is
impossible. Thus the exact order is \(n\), also for \(n=1\)
where the group is trivial. The finite Galois degree formula
gives \(|\operatorname{Gal}(k_n/k)|=n\), proving generation.
A positive power \(\Phi^d\) cannot be the identity on \(\Omega\):
its restriction to \(k_{d+1}\) has nonzero exponent modulo
\(d+1\). If \(m\mid n\), \(\Phi_n^m\) fixes \(k_m\), and in
the cyclic group of order \(n\) its order is \(n/m\). The
extension \(k_n/k_m\) is Galois with degree \(n/m\), by the
tower property and the degree formula. This proves the relative
assertion. See \cite[Proposition 4.20 and Example 7.17]{milne}
and \cite[Theorem 2.21]{ln}.
\end{proof}

\begin{proposition}[Restriction in exponent coordinates]\label{res:E05}
For every \(n\geq1\) the map
\[
     \lambda_n:(\mathbb Z/n\mathbb Z,+)
             \longrightarrow \operatorname{Gal}(k_n/k),
     \qquad \lambda_n(\bar a)=\Phi_n^a
\]
is a group isomorphism. For \(m\mid n\) let
\(\rho_{n,m}:\mathbb Z/n\mathbb Z\to\mathbb Z/m\mathbb Z\)
be reduction. Then
\[
          r_{k_n,k_m}\circ\lambda_n=\lambda_m\circ\rho_{n,m}.
\]
\end{proposition}
\begin{proof}
The group isomorphism follows from the generation and exact order
in Theorem~\ref{res:E04}. The reduction \(\rho_{n,m}\) is well-defined
since any multiple of \(n\) is a multiple of \(m\).
The automorphism \(\Phi_n\) preserves \(k_m\) and restricts to
\(\Phi_m\). Restriction is a homomorphism by Proposition~\ref{res:F07},
so it sends
\(\Phi_n^a\) to \(\Phi_m^a\) for every integer \(a\), including
negative exponents. This proves the displayed equality on every
residue class. In particular the isomorphisms use the same
distinguished generator at every stage. Compare
\cite[Examples 7.17 and 7.26]{milne}.
\end{proof}

\begin{theorem}[The absolute Galois group of a finite field]
\label{res:E06}
Endow
\[
    \widehat{\mathbb Z}
      =\varprojlim_{n}\mathbb Z/n\mathbb Z
\]
with its inverse-limit topology. The index is the set of positive
integers ordered by divisibility. The topology is that of compatible
residue families in the product of the finite discrete groups.
There is a topological group isomorphism
\[
         \Theta:\operatorname{Gal}(\Omega/k)
                          \xrightarrow{\ \sim\ }(\widehat{\mathbb Z},+)
\]
defined by \(\Theta(\sigma)_n=\lambda_n^{-1}(\sigma|_{k_n})\).
It sends \(\Phi\) to the family \(1=(1\bmod n)_n\).
\end{theorem}
\begin{proof}
By Corollary~\ref{res:E03}, the \(k_n\) are exactly the finite
Galois stages of the extension, with order given by
Proposition~\ref{res:E02}. The topological inverse-limit comparison
in Theorem~\ref{res:T06} identifies \(\operatorname{Gal}(\Omega/k)\)
with their compatible systems of automorphisms. The maps
\(\lambda_n^{-1}\) carry compatible systems to compatible residues
by Proposition~\ref{res:E05}. Conversely the \(\lambda_n\) give the
inverse map on compatible residues. Both maps between these limit
spaces are continuous, since every coordinate is given by a continuous
map between finite discrete spaces. The formula for \(\Phi\) follows
from \(\lambda_n(1)=\Phi_n\) for all \(n\).
This is the residue model of the profinite group in
\cite[Example 7.17]{milne} and \cite[Example 2.11.2]{rz}.
\end{proof}

\section{Intermediate fields in the residue model}

\begin{proposition}[A dense proper subgroup]\label{res:E07}
The ordinary cyclic subgroup \(\langle\Phi\rangle\) is dense and
proper in \(\operatorname{Gal}(\Omega/k)\). Its fixed field is \(k\).
\end{proposition}
\begin{proof}
Under \(\Theta\) its image is the group of integer residue families
\((a\bmod n)_n\) with \(a\in\mathbb Z\). Prescribing finitely many
coordinates of a compatible family is achieved by taking an integer
representative at a common multiple of those moduli. Hence this image
is dense.

To see that the image is proper, for each \(n=2^s d\) with \(d\) odd
choose the unique residue class \(b_n\bmod n\) with
\[
                      b_n\equiv0\pmod{2^s},
                \qquad b_n\equiv1\pmod d .
\]
This exists by the Chinese remainder theorem, including moduli
equal to \(1\). For \(m\mid n\), reduction of \(b_n\) satisfies the
corresponding two conditions for \(m\), so it equals \(b_m\).
The family \((b_n)_n\) is compatible. An integer representing it
would be divisible by all powers of \(2\), hence zero, but its
residue modulo \(3\) would be \(1\). Thus it does not come from
an integer. Finally the fixed elements of \(\Phi\) are the roots
of \(X^q-X\), namely \(k\) by Proposition~\ref{res:E01}.
Equivalently a subgroup and its closure have the same fixed field,
as in Theorem~\ref{res:C04}. See \cite[Example 7.17]{milne}
and \cite[Exercise 2.11.4]{rz}.
\end{proof}

\begin{theorem}[Open subgroups and finite intermediate fields]
\label{res:E08}
For each \(n\geq1\), write \(\pi_n:\widehat{\mathbb Z}\to
\mathbb Z/n\mathbb Z\) for projection. Then
\[
                   \ker\pi_n=n\widehat{\mathbb Z}.
\]
Under \(\Theta\) this subgroup corresponds to
\(\operatorname{Gal}(\Omega/k_n)\), has index \(n\), and has
fixed field \(k_n\). All open subgroups arise uniquely this way.
\end{theorem}
\begin{proof}
The image of the integers under \(\pi_n\) is the whole finite
coordinate group, so \(\pi_n\) is surjective and its kernel is
closed and open. It contains \(n\mathbb Z\), viewed in
\(\widehat{\mathbb Z}\) through the residue map. If \(a\in\ker\pi_n\)
and a neighbourhood of \(a\) specifies finitely many moduli,
choose \(N\) divisible by them and by \(n\). An integer
representative of \(a_N\) is a multiple of \(n\) since \(a_n=0\),
and realizes all the prescribed coordinates. Thus
\(\overline{n\mathbb Z}=\ker\pi_n\).

The group \(\widehat{\mathbb Z}\) is compact by the inverse-limit
description in Theorem~\ref{res:E06} and the compactness of
the product of finite discrete groups. Multiplication by \(n\)
is continuous, so its image is compact and thus closed in this
Hausdorff group. Its image contains \(n\mathbb Z\). Conversely, by
the density of \(\mathbb Z\) proved in Proposition~\ref{res:E07},
the image of every point under multiplication by \(n\) belongs
to \(\overline{n\mathbb Z}\). This proves the kernel formula.

Under \(\Theta\), the map \(\pi_n\) reads the restriction to
\(k_n\), expressed in exponent coordinates. Its kernel is
therefore the subgroup fixing \(k_n\), and the index and fixed
field follow from surjectivity and Theorem~\ref{res:C04}.
Every open subgroup has a finite intermediate fixed field by
Theorem~\ref{res:C06}. Corollary~\ref{res:E03} shows this
field is some \(k_n\), uniquely determined by its degree over \(k\).
The closed-subgroup correspondence recovers the subgroup.
Compare \cite[Theorem 7.13(b) and Example 7.17]{milne}
and \cite[Theorem 2.7.2]{rz}.
\end{proof}

\begin{example}[A closed subgroup with infinite index]\label{res:E09}
For a prime \(\ell\) the intermediate field
\[
              M_\ell=\bigcup_{r\geq0}k_{\ell^r}
\]
is infinite over \(k\) and is a proper subfield of \(\Omega\).
Its fixing subgroup corresponds to the closed nonopen subgroup
\[
                       \bigcap_{r\geq0}\ell^r\widehat{\mathbb Z}.
\]
\end{example}
\begin{proof}
The \(k_{\ell^r}\) form an increasing chain by Proposition~\ref{res:E02},
so their union is a field. Its finite subextensions have unbounded
degrees \(\ell^r\), making \([M_\ell:k]\) infinite. Take any prime
\(s\ne\ell\). If \(k_s\subseteq M_\ell\), a finite \(k\)-basis of
\(k_s\) would be contained in one \(k_{\ell^r}\) because the chain
is increasing. This would give \(s\mid\ell^r\), a contradiction.
The subgroup fixing the union is the intersection of the subgroups
fixing each layer. Theorem~\ref{res:E08} gives the displayed
intersection, which is closed. It is nonopen by the infinite degree
and Theorem~\ref{res:C06}, so the same theorem gives infinite index.
Compare
\cite[Remark 7.14 and Example 7.17]{milne}.
\end{proof}
